\section{Pruebas y CI/CD}
\label{sec:pruebas-cicd}

\subsection{Pirámide de pruebas}

La estrategia de pruebas cubre las cuatro capas clásicas de la pirámide~\cite{ammann2016testing},
más una capa de carga, cada una respondiendo una pregunta que las demás no pueden:

\begin{table}[H]
\centering
\small
\caption{Pirámide de pruebas del sistema, con resultados reales.}
\label{tab:piramide}
\begin{tabular}{@{}p{2.3cm}p{3.4cm}p{4.2cm}p{3.2cm}@{}}
\toprule
\textbf{Capa} & \textbf{Herramienta} & \textbf{Pregunta que responde} & \textbf{Resultado} \\
\midrule
Unitaria & JUnit 5 / Vitest / JUnit Android & ¿Esta clase, aislada, hace lo que dice? & 298/298 backend (5 servicios Java + Node + Python), 128/128 web, 33/33 móvil \\
Integración & Testcontainers & ¿Funciona contra un CockroachDB real, no un \emph{mock}? & \texttt{TicketRepositoryIntegrationTest} \\
Contrato & Pact (consumidor + proveedor) & ¿El backend sigue la forma exacta que el cliente espera? & 2/2 consumidor, 2/2 proveedor \\
E2E & Playwright / Compose Testing & ¿Todo el camino, navegador real incluido, funciona? & 9/9 en Chromium+Firefox+WebKit \\
Carga & Locust & ¿Aguanta concurrencia real? & Sección~\ref{sec:iso25010} \\
\bottomrule
\end{tabular}
\end{table}

Cada cifra de la fila «Unitaria» se obtiene ejecutando la suite correspondiente, no contándola a
mano: \texttt{mvn test} dentro de \texttt{services/auth-service}, \texttt{services/svc-principal},
\texttt{services/report-service}, \texttt{services/api-gateway} y \texttt{services/telemetry-service}
para los cinco servicios Java (62, 149, 37, 8 y 15 pruebas respectivamente); \texttt{npm test} (que
invoca \texttt{node --test}) dentro de \texttt{services/notification-service} (11); \texttt{pytest}
dentro de \texttt{services/ai-service} (16) — total backend, 298; \texttt{npm test} (que invoca
\texttt{vitest run}) dentro de \texttt{apps/web} (128); y \texttt{./gradlew
:app:testDebugUnitTest} dentro de \texttt{apps/mobile} para la capa unitaria móvil (33). Son los
mismos comandos, sin variación, que ejecutan los trabajos \texttt{test-backend}, \texttt{test-web}
y \texttt{test-mobile} de la Tabla~\ref{tab:cicd}, así que la cifra publicada es la que produce el
propio flujo automático en cada envío, no una recontada aparte.

\texttt{api-gateway} y \texttt{telemetry-service} tenían Dockerfile y pruebas JUnit reales en el
repositorio, pero ningún trabajo del flujo los ejecutaba ni los publicaba (hallazgo documentado
en el commit \texttt{2c5bba5}, fuera de alcance de esa corrección en su momento). De las 8 pruebas
nuevas de \texttt{api-gateway}, 7 son de \texttt{HttpMetricsGlobalFilter}: el único filtro del
\emph{gateway} sin ejercitar, cuya lógica reduce identificadores UUID en la ruta a un marcador de
posición (\texttt{/api/v1/tickets/\{id\}}) para no romper la cardinalidad de las métricas de
Prometheus.

\texttt{AuthGatewayFilter} tampoco tenía ninguna prueba, a pesar de ser el único punto que exige
autenticación sobre \texttt{/api/v1/tickets/**}: valida el token delegando en \texttt{auth-service}
por HTTP en vez de verificar la firma localmente (Sección~\ref{sec:contrato-openapi}). Sus 12
pruebas nuevas corren contra un doble real de \texttt{auth-service} (un
\texttt{com.sun.net.httpserver.HttpServer} local, no un \emph{mock} de la biblioteca HTTP),
incluida la caída completa del servicio: con el doble apagado a mitad de la prueba, el filtro
falla cerrado (401), nunca deja pasar. \texttt{HttpMetricsFilter} —la misma clase, copiada sin
variación en \texttt{auth-service}, \texttt{svc-principal} y \texttt{report-service}, ver el
comentario de la propia clase— tampoco tenía pruebas en ninguna de sus tres copias; las 12 nuevas
(4 por servicio) documentan además un hallazgo real: a diferencia de
\texttt{HttpMetricsGlobalFilter} en el \emph{gateway}, que normaliza cualquier UUID en la ruta sin
condición, esta clase solo normaliza cuando Spring MVC ya resolvió un \emph{handler} — una
petición que no matchea ninguna ruta cae a la URI cruda sin normalizar, la misma cardinalidad sin
límite que la clase existe para evitar. Se deja caracterizado, no corregido, para que quien lo
arregle tenga la prueba lista para actualizar junto con el fix.

\texttt{JwtAuthenticationFilter}, en \texttt{auth-service}, tampoco tenía prueba propia: la clase
que puebla el \texttt{SecurityContext} a partir de un token —la pieza que hace que
\texttt{@PreAuthorize("hasRole('ADMIN')")} funcione en \texttt{AdminUserController}— es distinta de
\texttt{JwtService} (la emisión y validación del JWT en sí, que ya tenía prueba propia desde antes).
Sus 7 pruebas nuevas usan un \texttt{JwtService} real, no simulado, e incluyen el caso que este
filtro más necesita rechazar: un token con la forma correcta pero firmado con un secreto distinto
al real. \texttt{report-service} tiene su propia copia de \texttt{AuthGatewayFilter} —no idéntica a
la de \texttt{ticket-service}: agrega una restricción propia (solo el rol \texttt{ADMIN} puede
consultar reportes, aunque el token sea válido) que no existe en ningún otro servicio, así que
necesitaba su propia prueba, no la del hermano de \texttt{svc-principal}. Sus 8 pruebas nuevas
confirman que un \texttt{TECNICO} o \texttt{CLIENTE} con sesión real reciben 403 (prohibido), no
401 (no autenticado) —son errores distintos que un cliente del API necesita distinguir.

\texttt{AdminBootstrap} (\texttt{auth-service}) crea la cuenta \texttt{ADMIN} por defecto en el
primer arranque, únicamente si todavía no existe ninguna —el guardia exacto que evita crear un
\texttt{ADMIN} duplicado en cada reinicio— y no tenía prueba. Sus 2 pruebas nuevas cubren ambas
ramas: sin \texttt{ADMIN} previo, se crea uno con la contraseña cifrada (nunca el texto plano de
configuración); con uno ya existente, no se guarda nada. \texttt{ReportQueryService.filtered()}
—la única lógica propia de ese servicio, el resto son delegaciones directas al puerto— filtra en
memoria por hasta tres criterios independientes, opcionales e insensibles a mayúsculas, y tampoco
tenía prueba; sus 8 pruebas nuevas cubren cada filtro por separado, la combinación de los tres a
la vez, y el caso que un filtro aplicado por error con \texttt{OR} en vez de \texttt{AND} dejaría
pasar sin que nadie lo notara.

\texttt{TicketListViewModel} (\texttt{apps/mobile}) tampoco tenía prueba: es el único
\emph{ViewModel} que combina dos filtros independientes —búsqueda por texto y estado— sobre el
mismo flujo, con una distinción sutil que solo una prueba puede proteger de verdad: el total y el
conteo de incumplimientos de SLA se calculan sobre \emph{todos} los tickets, no sobre los ya
filtrados, así que el resumen sigue siendo correcto aunque la lista visible esté filtrada. Sus 6
pruebas nuevas usan Turbine para suscribirse al \texttt{StateFlow} —a diferencia de
\texttt{TicketDetailViewModel} (un \texttt{MutableStateFlow} simple), este expone su estado vía
\texttt{combine(...).stateIn(..., WhileSubscribed(5\_000), ...)}, que solo empieza a producir
valores cuando alguien lo está observando— y cubren, además de la distinción de conteos, que la
búsqueda no distingue mayúsculas, que los dos filtros se combinan con \texttt{AND}, que tocar el
mismo filtro de estado dos veces lo quita, y que un fallo de red al refrescar conserva los
tickets ya guardados en caché en vez de vaciar la pantalla.

Cuatro piezas más quedaron sin prueba fuera del backend Java. \texttt{kafka\_consumer.py}
(\texttt{ai-service}): \texttt{classify()} ya tenía prueba propia, pero no
\texttt{\_handle\_ticket\_created()} (el guardado en Mongo y la publicación de
\texttt{ticket.classified}) ni \texttt{\_consume\_loop()} —sus 4 pruebas nuevas confirman,
entre otras cosas, que un mensaje de Kafka malformado no tumba el hilo consumidor para siempre,
tal como promete el comentario del propio archivo. \texttt{kafkaConsumer.js}
(\texttt{notification-service}): \texttt{dispatch()} ya tenía prueba, pero no
\texttt{handleMessage()} (el mapeo a documentos de Mongo); sus 3 pruebas nuevas confirman que un
tópico sin notificaciones que despachar ni siquiera abre la colección. \texttt{apps/web}: el
interceptor que adjunta \texttt{Authorization: Bearer} a cada petición (\texttt{apiClient.ts})
nunca se había probado en aislamiento —\texttt{authApi.test.ts} simula \texttt{axios.post}
entero, sin pasar por la cadena real de interceptores—; sus 3 pruebas nuevas invocan el
interceptor registrado directamente. Y \texttt{useAuth} (el \emph{hook} de autenticación) se
usaba de paso en \texttt{ProtectedRoute.test.tsx} pero nunca se ejercitaban
\texttt{login()}/\texttt{logout()} ni los estados \texttt{loading}/\texttt{error}; sus 7 pruebas
nuevas cubren el flujo completo, incluido un rechazo que no es una instancia de \texttt{Error}
(la rama que cae al mensaje genérico).

Dos piezas más, una en cada cliente, seguían sin prueba propia. \texttt{TicketDetailModal}
(\texttt{apps/web}): pide de nuevo el ticket al backend en vez de reusar el objeto de la fila
—para no mostrar un estado desactualizado si \texttt{EscalationScheduler} lo cambió entre que
se cargó la lista y el clic— y tiene una guardia de cierre (clic en el fondo cierra el modal,
clic dentro del contenido no) fácil de romper sin que ninguna prueba lo note; sus 9 pruebas
nuevas cubren el estado de carga, el de error, y confirman explícitamente que un clic dentro
del contenido \emph{no} dispara \texttt{onClose}. \texttt{TicketRepository}
(\texttt{apps/mobile}): la pieza con más lógica de negocio propia de toda la app móvil —el
reintento de \texttt{closeOnSite} (solo ante fallo de \emph{red}, nunca ante un 4xx/5xx real) y
el patrón caché-primero de \texttt{getTicket}— solo se ejercitaba indirectamente a través de
\texttt{TicketDetailViewModelTest}, que mockea el repositorio entero y nunca corre su lógica
real. Sus 9 pruebas nuevas mockean \texttt{android.util.Base64} de forma estática (no disponible
en pruebas JVM puras, razón por la que ningún otro archivo lo había tocado hasta ahora) y
confirman, entre otras cosas, que un error HTTP real no dispara reintento —a diferencia de un
fallo de red, que sí— y que agotar los 3 intentos sin conexión propaga la excepción real, no una
genérica.

Dos piezas más, pequeñas pero reales. El manejador de \texttt{GET /api/v1/notifications}
(\texttt{notification-service}) decide su filtro de Mongo según si llega \texttt{ticketId} por
query string o no —«si viene, filtra; si no, trae todo»— y no tenía prueba; sus 4 pruebas nuevas
invocan el manejador real extraído del router de Express (no hay \texttt{supertest} instalado en
este proyecto), sin levantar un servidor HTTP, y confirman también que el campo \texttt{\_id}
crudo de Mongo nunca sale en la respuesta. \texttt{SessionMetrics} (\texttt{auth-service}) registra
el \emph{Gauge} \texttt{app\_active\_sessions} exigido por el Módulo F de la Guía de Entrega 4,
pero nada confirmaba que quedara registrado con ese nombre exacto ni que reflejara el conteo real
del repositorio; sus 3 pruebas nuevas usan un \texttt{SimpleMeterRegistry} real y comprueban,
entre otras cosas, que el \emph{Gauge} se recalcula en cada lectura y no queda cacheado con el
primer valor.

Dos piezas más en \texttt{svc-principal}. \texttt{RepositoryTimingAspect} —el único punto de
instrumentación del histograma \texttt{crdb\_query\_duration\_seconds}, vía AOP sobre todo el
paquete de persistencia— no tenía prueba, y su garantía central (el cronómetro vive en un bloque
\texttt{finally}) nunca se había verificado: sus 4 pruebas nuevas confirman que una consulta que
falla igual cuenta para el histograma, y que la excepción real sigue propagándose sin que el
aspecto la trague. \texttt{TicketMapper} —la traducción campo a campo entre el dominio puro y su
entidad JPA— tampoco tenía prueba, pese a que un mapeo así es exactamente el tipo de código que
se rompe en silencio cuando se agrega un campo a un lado y se olvida el otro: ya pasó de verdad
con los campos de evidencia (\texttt{evidencePhoto}/\texttt{evidenceLatitude}/%
\texttt{evidenceLongitude}, Entregable 10), agregados a \texttt{Ticket} y \texttt{TicketJpaEntity}
por separado. Sus 4 pruebas nuevas incluyen un viaje de ida y vuelta campo por campo con
\texttt{usingRecursiveComparison()}, que fallaría de inmediato si un campo nuevo se agregara a un
lado sin el otro.

\texttt{TicketController} no tenía prueba propia —sus manejadores sí, extensamente—, pero
tiene una pieza de lógica que le pertenece solo a él: decodificar en Base64 la evidencia
fotográfica del cierre en sitio. Sus 7 pruebas nuevas (manejadores mockeados, sin \emph{MockMvc}
ni contexto de Spring) documentan un hallazgo real, no solo lo prueban: un
\texttt{evidencePhotoBase64} malformado —un error del lado del cliente móvil, no
necesariamente malicioso— no tiene manejo explícito y cae al
\texttt{@ExceptionHandler(Exception.class)} genérico de \texttt{GlobalExceptionHandler}, que
responde 500 «Error interno» en vez de un 400 claro. Se deja caracterizado, no corregido.

El único \emph{endpoint} HTTP de \texttt{ai-service} —\texttt{GET
/api/v1/ai/classifications/\{ticketId\}}, la contraparte de lectura de lo que
\texttt{\_handle\_ticket\_created()} escribe en Mongo— tampoco tenía prueba propia. Sus 4 pruebas
nuevas (llamando la función directo, sin \texttt{TestClient}/\texttt{httpx}, que ni siquiera está
en \texttt{requirements.txt}) confirman que un ticket sin clasificación todavía responde 404 con
el identificador en el mensaje —en vez de un 500 o, peor, un 200 con cuerpo vacío— y que la
proyección \texttt{\{"\_id": 0\}} realmente se pasa a Mongo, no solo se intenta: un
\texttt{ObjectId} crudo filtrado al cliente rompería la serialización JSON de FastAPI en cuanto
existiera al menos una clasificación guardada.

\texttt{AuthController}, en cambio, no tenía ninguna prueba propia pese a exponer los cinco
\emph{endpoints} de sesión del sistema. Sus 7 pruebas nuevas (con \texttt{AuthService}
mockeado, sin \emph{MockMvc} ni contexto de Spring) verifican que \texttt{login} extrae
\texttt{email}/\texttt{password} del \emph{request} en vez de reenviarlo entero, y sobre todo
cubren la única lógica que le pertenece al controlador y no al servicio: \texttt{validate()}
parsea el encabezado \texttt{Authorization} a mano (en vez de dejarlo a Spring Security, ver
comentario en la clase) para que un \emph{token} ausente o mal formado devuelva un 401 con el
envoltorio \texttt{ApiResponse} del dominio en lugar del manejador genérico del framework; un
encabezado nulo y uno sin el prefijo \texttt{"Bearer "} quedan probados por separado.

Las tres estrategias de correlación de incidencias —c0/c1/c2, el componente exacto que el
Escenario 4 del protocolo experimental (dos averías simultáneas en la misma zona) pone a prueba
(ver ADR-0008)— ya tenían una suite propia (\texttt{CorrelationStrategyTest}) desde que se
introdujeron, pero con dos huecos concretos. Ninguna prueba de c1 confirmaba que los tickets que
ya estaban en una incidencia sobrevivieran al unirse uno nuevo —solo se verificaba que el
recién llegado entrara, no que los anteriores no se perdieran, un bug de reemplazo-en-vez-de-agregar
habría pasado inadvertido—; y ninguna prueba de c2 verificaba el valor real que la conversión
minutos-a-segundos de la ventana le pasa al puerto de telemetría, porque los dobles de
\texttt{TelemetryQueryPort} usados ignoraban ese parámetro por completo. Sus 2 pruebas nuevas
cierran ambos huecos: una confirma que unirse a una candidata agrega el ticket sin descartar los
que ya estaban, y la otra captura el argumento real recibido por el puerto y confirma que una
ventana de 15 minutos configurada llega como 900 segundos, ni un valor crudo sin convertir ni
otra escala.

Dos piezas más, una en cada extremo del sistema. \texttt{TicketSummaryMapper}
(\texttt{report-service}) —el mismo tipo de traducción campo a campo entre dominio puro y
entidad JPA que ya causó un bug real en \texttt{ticket-service} (ver \texttt{TicketMapper},
Entregable 10)— tampoco tenía prueba, con un agravante propio: \texttt{TicketSummary} solo se
reconstruye a partir de eventos de Kafka (ver \texttt{ReportEventListener}), así que un campo
perdido en el mapeo no se nota en una escritura directa a la base, se nota semanas después, como
un reporte con un dato en blanco que nadie relaciona con el mapeo. Sus 4 pruebas nuevas incluyen
el mismo viaje de ida y vuelta con \texttt{usingRecursiveComparison()} que ya protege a
\texttt{TicketMapper}. \texttt{AuthRepository} (\texttt{apps/mobile}) tampoco tenía prueba: es
el único punto que decide si un login exitoso en el servidor realmente cuenta como sesión
iniciada en el cliente, envolviendo la llamada en \texttt{runCatching} y exigiendo
\texttt{data} no nulo con \texttt{requireNotNull} antes de guardar los tokens. Sus 5 pruebas
nuevas documentan, entre otras cosas, que una respuesta HTTP 200 con \texttt{data} nulo —un
caso real y no solo teórico, dado que \texttt{ApiResponse.data} es nullable en el contrato—
falla en vez de guardar tokens vacíos y reportar una sesión iniciada que en realidad no existe.

Dos piezas más, ambas en el extremo de PE-U1 que alimenta la correlación de incidencias.
\texttt{TelemetryStore} (\texttt{telemetry-service}) —el \emph{buffer} en memoria del que lee
\texttt{ZonaVentanaTelemetriaStrategy} (c2, ver Ronda 12 más arriba) para decidir si hay
evidencia real de avería antes de agrupar tickets— tampoco tenía prueba propia: las pruebas de
c2 mockean \texttt{TelemetryQueryPort} por completo y nunca ejercitan este buffer real. Sus 7
pruebas nuevas confirman, entre otras cosas, la garantía que el propio comentario de la clase
remarca —los eventos se devuelven ordenados por \emph{timestamp} de Lamport (orden causal), no
por orden de llegada— registrando eventos deliberadamente fuera de ese orden para que una
implementación que solo devolviera la lista tal cual, o la ordenara por hora de recepción, no
pasara la prueba; también confirman que \texttt{purgar()} descarta eventos vencidos en todas las
zonas sin tocar los vigentes. \texttt{PermissionCatalog} (\texttt{auth-service}) —el mapa
rol-a-permisos que se embebe como \emph{claim} en el \emph{access token} para que
\texttt{ticket-service}/\texttt{api-gateway} autoricen sin volver a consultar
\texttt{auth-service} en cada petición— tampoco tenía prueba; un permiso mal listado aquí se
propagaría en silencio a cada token emitido. Sus 4 pruebas nuevas incluyen una que compara los
tres roles entre sí, no solo contra el valor esperado, para atrapar el caso de que dos roles
compartieran por error la misma entrada del mapa.

Dos piezas más para cerrar esta ronda. \texttt{RefreshTokenMapper} (\texttt{auth-service}) —el
mismo tipo de mapeo campo a campo entre dominio puro y entidad JPA que ya causó un bug real en
\texttt{ticket-service} (ver \texttt{TicketMapper}, Entregable 10)— tampoco tenía prueba, con un
riesgo de seguridad concreto y no solo de datos en blanco: \texttt{revoked} es el único mecanismo
que impide reusar un \emph{refresh token} ya rotado o robado (ver
\texttt{TokenReuseDetectedException}); si este campo se perdiera en el mapeo, un token revocado
se recargaría desde la base como si nunca lo hubiera sido. Sus 5 pruebas nuevas incluyen un caso
dedicado a ese escenario, no solo el viaje de ida y vuelta genérico. Las funciones de mapeo de
\texttt{apps/mobile} —\texttt{colorForStatus}/\texttt{labelForStatus}/\texttt{colorForPriority}/
\texttt{labelForCategory} (\texttt{ui/theme/StatusColors.kt})— tampoco tenían prueba: son
\texttt{when} exhaustivos (el compilador de Kotlin exige cubrir todos los casos, así que nunca
falta una rama), pero el compilador no verifica que el \emph{valor} de cada rama sea el correcto
—un color copiado por error de un estado a otro, o una tilde mal escrita en una etiqueta, pasarían
el \emph{build} sin que nada lo note. Sus 8 pruebas nuevas confirman, entre otras cosas, que
ningún par de estados comparte color, y que \texttt{CRITICO} reutiliza a propósito el mismo tono
de \texttt{ESCALADO} (documentado explícitamente en el comentario de la clase), no por accidente.

\texttt{AuthService} (\texttt{auth-service}) ya tenía una suite propia de 10 pruebas —incluida la
detección de reuso de \emph{refresh token} con revocación masiva—, pero con seis ramas reales sin
ejercitar. Ninguna prueba cubría un correo ya registrado (\texttt{DuplicateEmailException}) ni un
\texttt{login} contra un correo que directamente no existe —distinto del caso ya probado de
contraseña incorrecta—, ambos debiendo fallar con el mismo \texttt{InvalidCredentialsException}
genérico para no revelar si un correo está registrado. Tampoco había prueba de \texttt{refresh}
con un \emph{hash} de token desconocido ni con uno vencido pero no revocado —dos causas de fallo
distintas del reuso ya cubierto—, ni de \texttt{logout} en absoluto: ni la revocación de una
sesión activa, ni su idempotencia (cerrar sesión dos veces, o después de un \texttt{refresh} que
ya revocó el mismo token, no debe volver a escribir en la base), ni el caso de un \emph{token}
desconocido. Tampoco \texttt{validate()} —el único método que traduce los \emph{claims} ya
verificados del JWT a la respuesta que consume el \emph{API Gateway}— tenía prueba propia. Sus 7
pruebas nuevas cierran los seis huecos.

\texttt{TicketQueryService} (\texttt{svc-principal}) —hallado revisando el reporte de cobertura
JaCoCo del módulo, no por inspección manual— ya tenía una suite de 9 pruebas, pero cubría menos
de la mitad de las ramas reales de la matriz de autorización por rol
(\texttt{TicketAuthorization}, ver su tabla \texttt{CLIENTE}/\texttt{TECNICO}/\texttt{ADMIN}).
De \texttt{getTicket()} solo estaba probado el camino \emph{prohibido} de \texttt{TECNICO} (otra
zona) y el de \texttt{ADMIN} contra un ticket inexistente —nunca el camino permitido de ninguno
de los dos roles con un ticket real. De \texttt{listTickets()} faltaban las tres ramas con filtro
de estado combinado con rol (\texttt{CLIENTE}, \texttt{TECNICO} y \texttt{ADMIN} con
\texttt{zone}+\texttt{status}), la rama de \texttt{ADMIN} con solo zona, y —la más importante— el
\emph{fallback} final para un rol no reconocido: sin prueba, nada garantizaba que un rol nuevo o
mal escrito cayera a lista vacía (\emph{fail-closed}) en vez de, por accidente de refactor,
colarse en alguna rama y filtrar tickets que no le correspondían. Sus 7 pruebas nuevas cierran
los siete huecos, incluida esa última como caso explícito.

Dos manejadores de comando más, encontrados leyendo directamente el reporte JaCoCo del módulo
(92.2\% y 95.1\% de cobertura respectivamente, no un 0\% como los hallazgos de rondas
anteriores). \texttt{AssignTechnicianHandler} tenía solo 2 pruebas —el camino prohibido de
\texttt{TECNICO} en otra zona, y la publicación del evento con \texttt{ADMIN}— sin cubrir ni el
ticket inexistente, ni \texttt{CLIENTE} (siempre prohibido), ni el camino \emph{permitido} de
\texttt{TECNICO} en su propia zona, ni \texttt{TECNICO} sin zona reconocible. Sus 4 pruebas
nuevas cierran esos cuatro huecos. \texttt{UpdateTicketStatusHandler} solo probaba
\texttt{slaBreached} en su rama verdadera (resuelto después del plazo); la rama negativa
—resuelto a tiempo— y el cortocircuito por \texttt{slaDeadline == null} (un ticket sin plazo
formal, por ejemplo antes de que \texttt{ai-service} lo clasifique) nunca se ejercitaban —sin
este segundo caso, quitar el chequeo de nulidad del \texttt{\&\&} habría tumbado con un
\texttt{NullPointerException} el cierre de cualquier ticket sin SLA asignado, sin que ninguna
prueba lo notara. Sus 2 pruebas nuevas cierran ambas ramas.

Tres huecos más, encontrados leyendo línea a línea el reporte HTML de JaCoCo (no solo el
porcentaje agregado por clase). \texttt{UpdateTicketStatusHandler} quedaba en 95.1\% incluso
después de la ronda anterior: la única línea marcada parcial era el \texttt{orElseThrow} del
ticket inexistente —ninguna prueba llamaba al manejador con un \texttt{id} que el repositorio no
encontrara. Su prueba nueva cierra esa última línea.
\texttt{TicketClassificationListener} —el lado receptor de la Saga por coreografía, que aplica
la clasificación de \texttt{ai-service} a un ticket— tenía tres pruebas, pero ninguna ejercitaba
\texttt{parseEnumOrNull()} más allá del caso feliz: ni un valor de categoría desconocido (el
escenario que el propio \emph{javadoc} del evento anticipa: \emph{"ai-service y ticket-service
pueden desincronizar su vocabulario"}), ni un payload que sencillamente no trae categoría ni
prioridad. Sus 2 pruebas nuevas confirman que ambos casos se ignoran en silencio —la categoría
queda \texttt{null}, el ticket se guarda igual— sin tumbar el \emph{listener}.
\texttt{TelemetryGrpcClientAdapter} ya tenía una prueba de regresión de un bug real (el plazo de
gRPC fijado una sola vez en el constructor en vez de por llamada), pero nunca ejercitaba su
propio bloque \texttt{catch} —el comportamiento \emph{fail-open} exigido por el ADR-0004: si
\texttt{telemetry-service} está caído, se trata como "sin evidencia", nunca se propaga la
excepción. Su prueba nueva apunta el adaptador a un puerto sin ningún servidor escuchando y
confirma que devuelve \texttt{false} en vez de lanzar.

\texttt{ThemeContext} (\texttt{apps/web}) ya tenía dos pruebas, pero ambas montaban el
\emph{provider} sin nada guardado en \texttt{localStorage}, así que siempre ejercitaban el
\emph{fallback} a la preferencia del sistema operativo (\texttt{matchMedia}, mockeada en
\texttt{test/setup.ts} sin preferencia oscura) —la rama real de "usuario que vuelve", con un
tema ya elegido en una visita anterior, nunca se probaba. Tampoco se confirmaba el propósito
mismo de guardar el tema en \texttt{localStorage}: la prueba original de \texttt{toggle} solo
verificaba el atributo \texttt{data-theme} del DOM, nunca que el nuevo valor quedara persistido.
Sus 2 pruebas nuevas cierran ambos huecos: una guarda \texttt{"dark"} en \texttt{localStorage}
antes de montar y confirma que el \emph{provider} lo lee sin consultar \texttt{matchMedia}, la
otra confirma que \texttt{toggle} deja el nuevo valor escrito en \texttt{localStorage}, no solo
reflejado en el DOM.

\texttt{CreateTicketModal} (\texttt{apps/web}, el único formulario que puede ver un
\texttt{CLIENTE}) ya tenía cuatro pruebas, pero desproporcionadamente pocas para sus 155 líneas:
cubrían el caso feliz, la validación con campos literalmente vacíos, y el error genérico de
\emph{fallo} de red, pero no el chequeo real —\texttt{title.trim()}, más estricto que
\texttt{!title}— con campos de puro espacio en blanco; tampoco se probaba el camino donde
teléfono/dirección sí se completan (recortados antes de enviarse, no solo el caso donde quedan
\texttt{undefined}), ni que cambiar la zona del \emph{select} realmente cambie la zona enviada
—las cuatro pruebas existentes dejaban siempre la zona por defecto—, ni la guardia de cierre
(clic en el fondo cierra, clic dentro del formulario no), el mismo patrón ya protegido en
\texttt{TicketDetailModal} (Ronda 6) pero nunca replicado aquí. Sus 5 pruebas nuevas cierran los
cinco huecos.

\texttt{TicketRowActions} (\texttt{apps/web}) ya tenía seis pruebas razonablemente completas
—estado, autoasignación de \texttt{TECNICO}, selector filtrado por zona de \texttt{ADMIN}, aviso
sin técnicos— pero con dos huecos reales. Ninguna probaba el \emph{guard}
\texttt{if (status === ticket.status) return}: sin él, reseleccionar en el \emph{select} el
mismo estado que el ticket ya tiene dispararía un \texttt{PATCH} redundante cada vez. Tampoco se
probaba, con rol \texttt{CLIENTE} (o sin rol), que el componente no ofreciera ningún control de
asignación —defensa en profundidad más allá de que \texttt{ConsolePage} ya lo oculte para ese
rol— ya que las seis pruebas existentes solo cubrían \texttt{TECNICO} y \texttt{ADMIN}. Sus 2
pruebas nuevas cierran ambos huecos.

\texttt{Badges} (\texttt{apps/web}) ya tenía cinco pruebas, pero con el mismo tipo de hueco que
\texttt{StatusColors.kt} en móvil (Ronda 15): \texttt{STATUS\_VAR} es un mapa manual estado-a-color
—a diferencia del \texttt{when} exhaustivo de Kotlin, aquí ni siquiera el compilador obliga a
cubrir los seis casos—, y las pruebas existentes solo verificaban la etiqueta traducida de
\emph{un} estado y \emph{una} prioridad, nunca el color en sí. Un color copiado por error de un
estado a otro habría pasado el \emph{build} sin que nada lo notara. \texttt{AssignedChip}
—reemplaza el UUID crudo del técnico por un chip genérico cuando lo ve un \texttt{CLIENTE},
simétrico a \texttt{UnassignedChip}, que sí tenía prueba— no tenía ninguna en absoluto. Sus 4
pruebas nuevas confirman que los seis estados tienen colores de fondo distintos entre sí, que
las cuatro prioridades tienen colores de borde distintos entre sí, que \texttt{CRITICO} reutiliza
a propósito el mismo tono que \texttt{ESCALADO} (documentado en el código, no accidental), y
cubren \texttt{AssignedChip} por primera vez.

\texttt{useTickets} (\texttt{apps/web}, el equivalente de consola a \texttt{TicketListViewModel}
en móvil, Ronda 5) ya tenía cinco pruebas, pero con cuatro ramas reales sin ejercitar. La
búsqueda también filtra por \texttt{ticketId} además de descripción —documentado explícitamente
en el propio código: sirve para verificar que una acción hecha desde otro cliente (la app móvil)
se propagó al mismo backend—, pero solo se probaba la búsqueda por descripción. Una búsqueda de
puro espacio en blanco nunca se probaba que se tratara como vacía. Las cuatro pruebas de filtro
existentes probaban búsqueda, zona y estado por separado, nunca combinados —donde un \texttt{OR}
por error en vez de \texttt{AND} pasaría desapercibido—. Y \texttt{refresh()} nunca se
ejercitaba: ninguna prueba confirmaba que realmente volviera a pedir los tickets al backend. Sus
4 pruebas nuevas cierran los cuatro huecos.

\texttt{AdminPage} (\texttt{apps/web}, alta de cuentas) ya tenía cinco pruebas, pero con cuatro
huecos reales para sus 215 líneas. Solo se probaba la creación exitosa dejando el rol por
defecto (\texttt{CLIENTE}, con \texttt{zone: null}); el camino real donde el rol elegido es
\texttt{TECNICO} —que exige una zona válida en el backend, ver
\texttt{AuthService.validateZoneForRole}— nunca enviaba la zona seleccionada en una prueba. El
único error de formulario probado era la validación local (campos vacíos); un fallo real de la
API (correo duplicado, por ejemplo) nunca se ejercitaba. Tampoco se confirmaba que el formulario
se limpiara después de crear una cuenta con éxito —solo el mensaje de éxito—, dejando abierta la
posibilidad de que datos de la cuenta anterior quedaran arrastrados a la siguiente. Y, igual que
en \texttt{CreateTicketModal} (Ronda 21), la validación real es \texttt{.trim()}, más estricta
que la única prueba existente de campos literalmente vacíos. Sus 4 pruebas nuevas cierran los
cuatro huecos.

\texttt{ReportsPage} (\texttt{apps/web}, lado de lectura del CQRS) tenía solo dos pruebas para
sus 98 líneas: el caso feliz con datos en los tres desgloses y el error de carga. Ningún
desglose vacío se probaba —\texttt{BreakdownCard} renderiza un guion en ese caso, rama que nunca
se ejercitaba—, tampoco el texto de carga antes de que la promesa resolviera, ni el cálculo real
del ancho de barra (\texttt{value/max*100\%}, con \texttt{max = Math.max(1, ...)} para evitar
dividir entre cero). El hallazgo más interesante: \texttt{ReportsPage} mantiene su \emph{propia}
copia de \texttt{STATUS\_VAR}, independiente de la de \texttt{Badges.tsx} (Ronda 23) —dos mapas
manuales que deben mantenerse sincronizados a mano, sin nada que los ate—, y su
\texttt{colorFor(key) ?? 'var(--color-navy)'} para un estado no reconocido nunca se probaba. Sus
4 pruebas nuevas cierran los cuatro huecos.

\texttt{LoginPage} (\texttt{apps/web}) tenía solo dos pruebas para sus 281 líneas —la mayoría
constantes de estilo, pero con lógica real de todas formas—: la validación local de campos
vacíos y que escribir la limpia. Ninguna de las dos ejercitaba \texttt{authApi.login} de verdad
—quedaba sin mockear, así que cualquier intento de \emph{submit} real habría fallado por red
antes de poder observarse—, así que el flujo completo de autenticación nunca se probaba de
punta a punta: ni que un \emph{login} exitoso navegue a \texttt{/main}, ni que el correo se
recorte antes de enviarse (la contraseña, deliberadamente, no), ni que un fallo del backend
muestre el mensaje de error real en vez de quedar mudo. Tampoco se probaba el botón del ojo que
alterna el campo de contraseña entre oculto y visible. Sus 4 pruebas nuevas —mockeando
\texttt{authApi.login} con el mismo criterio que \texttt{useAuth.test.tsx}— cierran los cuatro
huecos.

\texttt{AppLayout} (\texttt{apps/web}, cabecera y navegación de toda página autenticada) ya
tenía cuatro pruebas, pero con tres huecos reales. \texttt{CLIENTE} —el rol más común del
sistema— nunca se probaba de forma explícita para la regla de RBAC que oculta
\texttt{Administración}/\texttt{Reportes}: solo se cubrían "sin sesión" y \texttt{ADMIN}. El
botón \texttt{Salir} solo se confirmaba cerrando la sesión (\texttt{isLoggedIn() === false});
que también navegara a \texttt{/login} nunca se verificaba de punta a punta. Y, cuarta aparición
del mismo patrón de hueco (\texttt{Badges.tsx}, Ronda 23; \texttt{ReportsPage}, Ronda 26;
\texttt{StatusColors.kt} en móvil, Ronda 15): \texttt{ROLE\_BADGE\_COLOR} es otro mapa manual
rol-a-color, y las pruebas existentes solo verificaban la etiqueta traducida
(\texttt{"Administrador"}, \texttt{"Técnico"}), nunca que los tres roles tuvieran colores de
insignia distintos entre sí. Sus 3 pruebas nuevas cierran los tres huecos.

\texttt{SettingsPage} (\texttt{apps/web}) ya tenía dos pruebas —cambiar idioma, cambiar
tema—, pero ninguna cubría el \emph{guard} \texttt{mode !== 'light' \&\& toggle()}: sin él,
reclicar el botón del tema ya activo llamaría a \texttt{toggle()} igual y cambiaría al otro
tema por error, el mismo tipo de \emph{guard} redundante ya protegido en
\texttt{TicketRowActions} (Ronda 22) para el estado del ticket. Tampoco se comparaba
\texttt{pillStyle(active)} entre el botón activo y el inactivo —quinta aparición del patrón de
hueco de color/estilo nunca verificado (\texttt{Badges.tsx}, Ronda 23; \texttt{ReportsPage},
Ronda 26; \texttt{StatusColors.kt} en móvil, Ronda 15; \texttt{AppLayout}, Ronda 28)—. Sus 2
pruebas nuevas cierran ambos huecos.

Se cambió de capa para esta ronda: dos traductores centrales de excepciones a respuestas HTTP
(\texttt{@RestControllerAdvice}) sin ninguna prueba, a diferencia de su equivalente ya cubierto
en \texttt{svc-principal}. El de \texttt{auth-service} traduce seis excepciones distintas —cada
una a su código de estado correcto (404/400/401/403/500), incluida la única lógica real del
archivo: unir varios errores de campo de validación con \texttt{"; "} o caer a un mensaje
genérico cuando no hay ninguno—; sus 8 pruebas nuevas cubren los seis \emph{handlers}, ambas
ramas de la unión de errores, y confirman explícitamente que \texttt{handleAccessDenied} usa un
mensaje propio fijo en vez de reenviar \texttt{ex.getMessage()} —si alguien lo cambiara por
error, podría filtrar detalles internos de Spring Security en la respuesta—. El de
\texttt{report-service} es más simple (\texttt{ForbiddenException} → 403, cualquier otra cosa →
500); sus 2 pruebas nuevas lo cubren por completo.

\texttt{TicketRepositoryAdapter} era, según el reporte JaCoCo, la clase peor cubierta de todo el
módulo (29.2\%): no tenía ninguna prueba unitaria propia, solo se ejercitaba de forma indirecta
a través de \texttt{TicketRepositoryIntegrationTest} (Testcontainers, un CockroachDB real), que
no corre localmente en Windows por la limitación de \emph{Docker outside of Docker} ya
documentada en la Sección de amenazas —así que en la práctica, fuera del \emph{runner} de GitHub
Actions, esta clase no tenía cobertura real alguna. Sus 8 pruebas nuevas cubren los ocho métodos
del adaptador (siete consultas más \texttt{save()}), usando el \texttt{TicketMapper} real, no
mockeado, para que un cambio accidental en el mapeo también rompa esta prueba.
\texttt{IncidenciaRepositoryAdapter}, su equivalente para incidencias, ya tenía cinco pruebas
pero dejaba sin cubrir \texttt{findByZone()}; su prueba nueva cierra ese último método.
\texttt{TicketRepositoryAdapter} queda en 100\% de cobertura de línea según el propio reporte
JaCoCo tras esta ronda.

\subsubsection{Contrato de interfaz (OpenAPI)}
\label{sec:contrato-openapi}

Antes de que exista una sola prueba de contrato hace falta el contrato en sí. El archivo
\texttt{docs/api/openapi.yaml} (OpenAPI~3.0.3) describe, de forma verificable por herramienta y no
solo en prosa, las diecisiete operaciones REST de los cinco microservicios: registro, login,
refresco, logout, validación de sesión y administración de cuentas (\texttt{auth-service}, siete
operaciones); creación, listado, detalle, cambio de estado y asignación de tickets
(\texttt{ticket-service}, cinco); notificaciones (\texttt{notification-service}, una);
clasificación automática (\texttt{ai-service}, una); y los tres puntos de agregación de
\texttt{report-service}, incluida la exportación a CSV. Cada una de las diecisiete declara su
esquema de respuesta exitosa, los códigos de estado posibles, y el esquema de seguridad
(\texttt{bearerAuth}, JWT) cuando aplica; se valida automáticamente en el flujo con Spectral
(Sección~\ref{sec:validacion-contrato}), que rompe la construcción si alguna deja de cumplirlo.

Este archivo no es un anexo estático: es la fuente que el contrato Pact de la siguiente
subsección referencia como forma esperada, y es lo que permite a un integrante nuevo entender la
API completa sin leer el código de los cinco servicios. Se actualizó para esta entrega —seguía
describiendo los puertos directos de la Entrega 2 como si fueran la única vía de acceso, sin
mencionar el API Gateway que desde la Entrega 4 es el punto de entrada real de ambos clientes
(Sección~\ref{sec:arquitectura-e4})— para reflejar las dos rutas reales: el \emph{gateway} en el
puerto 8000 como camino de producción, y los puertos directos de cada servicio conservados solo
como vía de depuración local.

\subsubsection{Validación automática del contrato}
\label{sec:validacion-contrato}

Que el contrato declare un esquema no basta si nada impide que una operación nueva se añada sin
uno: la Entrega 4 se calificó con 4 operaciones sin esquema de respuesta señaladas por el
profesor, y al revisar las diecisiete resultaron ser once, no cuatro. Para que ese hueco no
vuelva a pasar desapercibido, el trabajo \texttt{validate-openapi} del flujo
(\texttt{.github/workflows/ci-cd.yml}) corre \texttt{Spectral} sobre
\texttt{docs/api/openapi.yaml} en cada envío y cada \emph{pull request}, con una regla propia
declarada en \texttt{.spectral.yaml} (\texttt{response-schema-required}) que exige
\texttt{content.application/json.schema} en toda respuesta 2xx; si una operación futura lo omite,
la construcción falla con el nombre exacto de la operación incumplida, no con una advertencia que
se pueda ignorar.

\subsubsection{Pruebas de contrato}

Se adoptó \emph{consumer-driven contract testing} con Pact entre \texttt{apps/web} (consumidor) y
\texttt{ticket-service} (proveedor). El consumidor genera un archivo de contrato
(\texttt{pacts/soporte-web-ticket-service.json}) a partir de las expectativas reales del cliente
sobre la forma de la respuesta; el proveedor lo verifica reproduciendo cada interacción contra el
servicio real. Una particularidad de esta arquitectura obligó a una decisión de diseño: como
\texttt{ticket-service} valida cada JWT llamando por HTTP a \texttt{auth-service}
(\texttt{AuthGatewayFilter}) en vez de verificar la firma localmente, la verificación del
proveedor no puede correr en un contexto de prueba aislado —necesita el \emph{stack} completo
levantado y un JWT real obtenido de un inicio de sesión real.

Un segundo ajuste, encontrado al ejecutar la verificación por primera vez, fue de modelado del
contrato: los campos \texttt{technicianId} y \texttt{resolvedAt} son legítimamente nulos o no
nulos según el estado real del ticket, y un solo ejemplo con matiz de tipo (\texttt{like(null)})
no puede expresar ambos casos a la vez sobre el mismo campo. Se adoptó el patrón oficial de Pact
para colecciones con formas mixtas (\texttt{arrayContaining} con dos variantes representativas),
en vez de forzar una condición más débil o más fuerte que la que el sistema realmente cumple.

\subsection{Pipeline de CI/CD}

El archivo \texttt{.github/workflows/ci-cd.yml} define siete \emph{jobs} en un grafo dirigido
acíclico, siguiendo la cultura DevOps de que ningún cambio llegue a un artefacto publicado sin
pasar por los mismos controles~\cite{kim2021devops,humble2010continuous}:

\begin{table}[H]
\centering
\small
\caption{Los siete jobs del pipeline.}
\label{tab:cicd}
\begin{tabular}{@{}p{0.6cm}p{2.6cm}p{9.5cm}@{}}
\toprule
\textbf{\#} & \textbf{Job} & \textbf{Qué hace} \\
\midrule
1 & \texttt{lint} & ESLint (web) + Android Lint (móvil) \\
2 & \texttt{test-backend} & Los cinco servicios Java con pruebas automatizadas + Node + Python, pruebas unitarias e integración \\
3 & \texttt{test-web} & Vitest, incluido el contrato Pact del consumidor \\
4 & \texttt{test-mobile} & Unitarias JVM + instrumentadas en un emulador Android headless real \\
5 & \texttt{build-images} & Construye y publica las ocho imágenes Docker en GHCR (solo en \emph{push}) \\
6 & \texttt{build-mobile-apk} & Compila el APK de depuración y lo publica como artefacto \\
7 & \texttt{integration} & Levanta el \emph{stack} completo y corre Pact proveedor + Playwright contra el backend real \\
\bottomrule
\end{tabular}
\end{table}

Los \emph{jobs} 5 y 6 declaran \texttt{needs} sobre los \emph{jobs} de prueba correspondientes:
GitHub Actions no los ejecuta hasta que sus dependencias terminen en verde. Esto no siempre fue
así de completo: en una revisión externa del pipeline se encontró que \texttt{build-images} solo
dependía de \texttt{test-backend} y \texttt{test-web}, así que publicaba imágenes en GHCR aunque
\texttt{lint} o \texttt{test-mobile} estuvieran en rojo -- no era la compuerta de calidad real que
se pretendía. Se corrigió agregando esas dos dependencias faltantes. Una segunda revisión
encontró que aún faltaba la más importante: \texttt{integration} -- el único \emph{job} que
prueba contra el \emph{stack} real y no contra \emph{mocks} -- corría en \emph{paralelo} con
\texttt{build-images} en vez de antes, así que una imagen podía publicarse aunque la integración
real fallara. Se agregó también a \texttt{needs}, sin crear una dependencia circular porque
\texttt{integration} construye sus propias imágenes en local (\texttt{docker compose up
--build}), no depende de las que publica GHCR. La declaración final es
\texttt{needs: [lint, test-backend, test-web, test-mobile, integration]}, para que ningún
artefacto se publique sin que los cinco jobs de verificación pasen primero.

\subsection{Depuración real del job \texttt{integration}}
\label{sec:cicd-depuracion}

Los primeros 6 \emph{jobs} pasaban de forma consistente desde su introducción. El séptimo,
\texttt{integration}, falló en \emph{todas} las corridas sobre GitHub Actions durante varios
días, pese a que el mismo \texttt{docker compose up -d --build} que ese \emph{job} ejecuta
funcionaba sin problema en las máquinas de desarrollo del equipo. Esa asimetría -- funciona en
local, falla siempre en el \emph{runner} -- fue la pista de que la causa era estructural, no una
falla puntual. El equipo instaló la CLI oficial de GitHub (\texttt{gh}) para acceder al log real
de cada paso (la API pública solo devuelve \enquote{\texttt{Process completed with exit code 1}},
sin el comando que realmente falló), y diagnosticó tres causas raíz independientes, cada una
confirmada con evidencia real antes de corregirla:

\begin{enumerate}[nosep]
  \item \textbf{Candado circular en el healthcheck de CockroachDB.} El healthcheck de los 3
        nodos usaba \texttt{/health?ready=1} -- que solo responde \texttt{OK} una vez que el
        nodo ya pertenece a un clúster \emph{inicializado} -- mientras que \texttt{db-init} (quien
        ejecuta \texttt{cockroach init}) esperaba a que ese mismo healthcheck pasara primero.
        Nadie quedaba listo porque nadie los inicializaba, y nadie los inicializaba porque nadie
        quedaba listo. En desarrollo esto no se manifestaba porque los volúmenes de Docker ya
        traían un clúster inicializado de sesiones anteriores; un \emph{runner} de GitHub
        Actions arranca siempre desde un volumen vacío, así que el candado se activaba sin
        excepción. Un primer intento de aumentar el presupuesto de tiempo del healthcheck
        \emph{no} resolvió el problema -- solo retrasó la falla de 70\,s a 3\,min 21\,s --,
        confirmando que la causa era estructural y no de tiempo. Se corrigió cambiando el
        healthcheck a \texttt{/health} (solo \emph{liveness}) y agregando reintentos reales
        alrededor de \texttt{cockroach init}.
  \item \textbf{Datos de demo del contrato Pact, nunca versionados.} El método \texttt{@State}
        de la verificación del proveedor asumía datos de demo ya cargados por un script que
        resultó no existir en ningún commit del repositorio -- solo había existido, sin
        documentarlo como tal, en el volumen de Docker de quien lo escribió originalmente. Se
        corrigió haciendo que el propio \texttt{@State} inserte esos datos por SQL directo, que
        es precisamente para lo que existe ese mecanismo de Pact.
  \item \textbf{Playwright solo instalaba Chromium.} \texttt{playwright.config.ts} define 3
        proyectos (Chromium, Firefox, WebKit -- requisito de compatibilidad de esta entrega),
        pero el paso de instalación del \emph{workflow} tenía el argumento
        \texttt{chromium} agregado, restringiendo la instalación a un solo navegador. Se
        corrigió quitando el argumento.
\end{enumerate}

Las tres correcciones se verificaron de punta a punta contra el \emph{stack} real -- nunca
simuladas -- antes de subirlas, y el pipeline completo quedó en verde sobre un mismo commit. El
patrón común a los tres problemas es el mismo: algo que \emph{parecía} funcionar porque la
máquina de desarrollo conservaba un estado acumulado (un volumen ya inicializado, datos cargados
a mano, un navegador ya instalado) que un entorno realmente limpio nunca reproduce por su cuenta.
El detalle completo de cada diagnóstico, incluidos los registros reales y los intentos que no
funcionaron, está en \texttt{docs/ci\_cd/informe\_ci\_cd.tex}.

\subsection{Cobertura y complejidad ciclomática}

La cobertura de \texttt{ticket-service}, medida con JaCoCo, no se resume en una sola cifra:
la Tabla~\ref{tab:cobertura-completa} reporta las tres métricas que la propia herramienta
entrega, cada una nombrada, en vez de citar solo la más favorable (líneas) sin decir que es la
de líneas.

\begin{table}[H]
\centering
\caption{Cobertura de \texttt{ticket-service}, las tres métricas de JaCoCo}
\label{tab:cobertura-completa}
\begin{tabular}{@{}lrrr@{}}
\toprule
\textbf{Métrica} & \textbf{Cubierto} & \textbf{Total} & \textbf{\%} \\
\midrule
Líneas & 555 & 684 & \textbf{81.1\,\%} \\
Instrucciones & 2582 & 3582 & 72.1\,\% \\
Ramas & 88 & 154 & 57.1\,\% \\
\bottomrule
\end{tabular}
\end{table}

Las tres cifras son ciertas simultáneamente y miden cosas distintas: líneas cuenta sentencias
ejecutadas, instrucciones cuenta a nivel de bytecode (más granular, penaliza más una rama sin
cubrir dentro de una misma línea), y ramas cuenta específicamente las bifurcaciones
condicionales (\texttt{if}/\texttt{switch}/operadores lógicos de corto circuito) — es la más
exigente y la que queda más lejos del objetivo del 70\,\% que el proyecto se propuso desde la
Entrega 3. El paquete concreto que más arrastra la cifra de ramas hacia abajo es
\texttt{infrastructure.security} (filtro y configuración de JWT), con \textbf{0 de 38 líneas
cubiertas}: ningún test ejercita las ramas de token inválido/expirado/ausente directamente sobre
esas clases —sí se ejercitan indirectamente por las pruebas de integración HTTP 401/403, pero
JaCoCo no atribuye esa cobertura a la clase que realmente decide el rechazo—, y queda declarado
aquí como una brecha real en vez de diluirse en el promedio general.

El reporte XML de JaCoCo agrega, a nivel de \texttt{<report>}, un total de 551/679 líneas —no
las 555/684 que resultan de sumar el CSV fila por fila. La diferencia no es un error de cálculo:
cinco archivos fuente (\texttt{AuthGatewayFilter.java}, \texttt{Incidencia.java},
\texttt{Ticket.java}, \texttt{IncidenciaJpaEntity.java} y \texttt{TicketJpaEntity.java}) contienen
dos clases cada uno —una externa y una interna o lambda generada por el compilador—, y el CSV
reporta la cobertura por clase, de modo que las líneas físicas compartidas entre ambas clases del
mismo archivo se cuentan dos veces al sumarlas. El XML, en cambio, agrega por archivo fuente y
cuenta cada línea física una sola vez. Las cifras de instrucciones y ramas no sufren este efecto.

Ninguna de las tres cifras incluye el código generado automáticamente por \texttt{protobuf}/
\texttt{grpc-java} a partir de \texttt{telemetry.proto} (\emph{getters}, \emph{setters} y
constructores mecánicos sin lógica propia, que ninguna convención de pruebas exige cubrir a
mano) — se excluye explícitamente en el \texttt{pom.xml}, con un comentario que documenta por
qué. Sin esa exclusión, ese código generado por sí solo diluía la cifra reportada de 59.3\,\%
(antes de que \texttt{telemetry-service}/\texttt{CORREL} existieran) a un engañoso 42.7\,\%,
contando en contra código que nadie prueba por convención en ningún proyecto Java real. El
reporte HTML de JaCoCo que respalda estas tres cifras —no solo una tabla escrita a mano— está
publicado en \texttt{docs/evidencias/jacoco-svc-principal/}, junto con el archivo crudo
\texttt{jacoco.csv} del que se recalcularon directamente (no se re-tipearon a mano) y el binario
crudo (\texttt{jacoco-svc-principal-raw.exec}) para quien quiera regenerarlo o verificarlo con
otra herramienta.

Hasta esta entrega, el objetivo del 70\,\% del Módulo A no se hacía cumplir: JaCoCo vivía en un
perfil Maven que no se activaba por defecto (\texttt{mvn test -Pcoverage}), el flujo de CI corría
\texttt{mvn test} sin ese perfil, y la ejecución \texttt{check} —la que compara contra un
umbral y puede romper la construcción— no existía en ningún sitio. El pipeline nunca medía nada,
pese a que el propio comentario del \texttt{pom.xml} decía que el reporte «es usado por el
pipeline». Se movió el complemento fuera del perfil (ahora siempre activo) y se agregó una
ejecución \texttt{check} con un umbral del 70\,\% sobre el contador de instrucciones, ligada a la
misma fase \texttt{test} que ya corre en \texttt{test-backend}: si la cobertura cae por debajo de
ese umbral, \texttt{mvn test} falla con la razón exacta.

Antes de escribir pruebas nuevas para esta entrega, el código propio (sin el generado) ya estaba
en 64.4\,\%. Se identificaron con el propio reporte de JaCoCo las clases reales con mayor número
de líneas sin cubrir — no una elección arbitraria — y se agregaron 31 pruebas unitarias nuevas
(87 en total en el módulo, todas en verde) sobre, entre otras: \texttt{CorrelationService} (el
orquestador de \texttt{CORREL}, que llegó a esta entrega prácticamente sin pruebas propias), los
tres observadores concretos del patrón Observer (\texttt{EscalationLoggingObserver},
\texttt{EscalationMetricsObserver}, \texttt{EscalationEventPublisherObserver} — cerrando un hueco
real en un patrón que el propio manuscrito documenta como implementado), \texttt{
GlobalExceptionHandler}, y el mapeo/adaptador de persistencia de \texttt{Incidencia}.

La complejidad ciclomática promedio del mismo módulo, medida con la herramienta \texttt{lizard},
es de \textbf{1.8}, muy por debajo del umbral de 10 — ninguna función individual supera siquiera
el umbral de alarma de 15, resultado consistente con el refactor a métodos pequeños de la
Sección~\ref{sec:hexagonal-e4}.
